\documentclass[aspectratio=169]{beamer}
\usepackage{tikz}
\usepackage{graphicx}
\usepackage{listings}
\usepackage{xcolor}
\usepackage{hyperref}
\usepackage[spanish,es-noquoting]{babel}
\usepackage[round]{natbib}

\lstset{
  basicstyle=\ttfamily\small,
  keywordstyle=\color{blue},
  stringstyle=\color{red},
  commentstyle=\color{gray},
  showstringspaces=false
}
% Minimal look
\usetheme{Warsaw}
%\setbeamertemplate{headline}
\setbeamertemplate{footline}{
  \leavevmode%
  \hbox{%
    % Left: Section name
    \begin{beamercolorbox}[wd=.33\paperwidth,ht=2.25ex,dp=1ex,left]{author in head/foot}%
      \hspace*{2ex}\insertsection
    \end{beamercolorbox}%
    % Center: Presentation title
    \begin{beamercolorbox}[wd=.34\paperwidth,ht=2.25ex,dp=1ex,center]{title in head/foot}%
      \insertshorttitle
    \end{beamercolorbox}%
    % Right: Slide numbers
    \begin{beamercolorbox}[wd=.33\paperwidth,ht=2.25ex,dp=1ex,right]{date in head/foot}%
      \insertframenumber{} / \inserttotalframenumber\hspace*{2ex}
    \end{beamercolorbox}%
  }%
  \vskip0pt%
}

\usecolortheme{default}
\usefonttheme{serif}
\graphicspath{{../}{./}}
\setbeamertemplate{caption}[numbered]

% Title info
\title[Módulo 3: Más allá del Texto]{Módulo 3: Más allá del Texto}
\subtitle{Análisis de Documentos y Creatividad}
\author{Ricardo Reyes-Jimenez}
\institute{IA para Todos}
\date{\today}



\begin{document}

% Title slide
\begin{frame}
  \titlepage
\end{frame}

% Reset background for normal slides
\usebackgroundtemplate{}



\begin{frame}[fragile]{Objetivo de la Sesión}


\textbf{Al finalizar este módulo, los participantes serán capaces de:}

1. Utilizar la IA para resumir y extraer información de documentos largos.
2. Explorar el uso de IA como un compañero para lluvia de ideas (Brainstorming).
3. Entender las limitaciones y consideraciones de privacidad básicas al usar estas herramientas.

\end{frame}



\begin{frame}[fragile]{1. Analizando Documentos Largos}


\textbf{La IA puede leer por nosotros, pero nosotros debemos validar.}

\begin{itemize}
  \item \textbf{Caso Práctico:} Un abogado (o cualquier persona firmando un contrato de alquiler) necesita revisar un documento de 10 páginas para encontrar cláusulas abusivas.
  \item \textbf{Uso de IA:} ``Adjunto el contrato de alquiler. Actúa como un revisor imparcial. Extrae las cláusulas relacionadas con penalizaciones por terminación anticipada y resúmelas en puntos breves. Señala si hay alguna condición inusual.''
\end{itemize}
\textbf{Precaución:} Nunca subas información personal, financiera confidencial o secretos comerciales a modelos públicos. Si necesitas hacerlo, usa versiones corporativas seguras.

\end{frame}



\begin{frame}[fragile]{2. El Compañero de Lluvia de Ideas}


\textbf{A veces, el mayor desafío es empezar (el síndrome de la hoja en blanco).}

\begin{itemize}
  \item \textbf{Caso Práctico:} Un emprendedor (o alguien organizando un evento) no sabe qué nombre ponerle a su nuevo proyecto de repostería.
  \item \textbf{Uso de IA:} ``Voy a abrir un negocio de repostería especializada en pasteles sin azúcar para diabéticos. La vibra es moderna y alegre. Dame 10 opciones de nombres creativos, en español, que no suenen médicos ni aburridos. Explica brevemente por qué sugieres cada uno.''
\end{itemize}
\textbf{Resultado:} Múltiples ideas en segundos para encender tu propia creatividad \cite{brynjolfsson2023}.

\end{frame}



\begin{frame}[fragile]{3. Privacidad y Alucinaciones}


\textbf{Dos reglas de oro para la vida diaria con IA:}

1. \textbf{La IA puede ``Alucinar'':} Los modelos pueden inventar datos con muchísima confianza. Siempre verifica la información importante (fechas, leyes, diagnósticos, datos históricos).
2. \textbf{Tu información entrena al modelo:} A menos que pagues por una versión privada o de empresa, asume que lo que escribes podría ser leído por otros. No compartas contraseñas, documentos de identidad ni secretos.

\end{frame}



\begin{frame}[fragile]{Recapitulando el Módulo}


\textbf{Hoy hemos aprendido:}
\begin{itemize}
  \item A usar la IA como un analista de documentos largos.
  \item Que la IA es excelente para superar bloqueos creativos y generar lluvia de ideas.
  \item A mantenernos seguros cuidando qué información compartimos y verificando las respuestas.
\end{itemize}
\textbf{¡Sigan aprendiendo, sigan trabajando y sigan disfrutando de la Inteligencia Artificial!}
\end{frame}




\begin{frame}[allowframebreaks]{References}
  \bibliographystyle{apalike}
  \bibliography{references}
\end{frame}


\end{document}